
\subsubsection{Problem Formulation}

Let D be a training corpus and B a benchmark dataset. A filtering strategy f: D $\rightarrow D_f$ selects a subset for training. Contamination detection c: D $\times$ B $\rightarrow$ {0, 1} labels each example as contaminated or clean.

Given a model M trained on $D_f$, let $\alpha _{i}^B$ denote the attribution score of training example i for benchmark performance. The Contamination Contribution Ratio is:

$CCR(M, D_f, B) = \Sigma _{i: c(d_i, B)=1} \alpha _{i}^B / \Sigma _i \alpha _{i}^B$

CCR measures the fraction of benchmark-attributed influence originating from contaminated examples.

\subsubsection{Contamination Detection}

Contamination is detected using 8-gram overlap between training documents and benchmark questions, following recommendations that n > 8 leads to false negatives while n < 8 produces excessive false positives.

\subsubsection{Attribution via TRAK}

Attribution scores are computed using TRAK, which approximates influence functions through random projection. TRAK expresses the counterfactual effect of removing example i on benchmark performance through projected gradient dot products.

\subsubsection{Amplification Index}

The Amplification Index (AI) measures differential contamination effects:

AI = $\Delta Acc_contaminated - \Delta Acc_clean$

where $\Delta Acc_contaminated$ is accuracy change on a potentially contaminated benchmark (MMLU) and $\Delta Acc_clean$ is accuracy change on a time-stratified clean benchmark (post-training cutoff). Positive AI indicates disproportionate improvement on contaminated benchmarks.

\subsubsection{Influence Fragility Ratio}

The Influence Fragility Ratio (IFR) characterizes replaceability:

IFR(i) = $\Delta Acc_mask(i) / \Delta Acc_retrain(i)$

IFR > 1 indicates the example's influence cannot be substituted by other training examples.
